% GENERATED FILE. Edit canonical lessons, then run npm run generate:books.
% canonical-source-hash: fnv1a64:a508a29128aee2f3
% canonical-lessons: RU-C102-boots, RU-C102-skirt, RU-C102-coat, RU-C102-shoes, RU-C102-socks

\chapter{Boots, Skirt, Coat, Shoes, Socks}
\label{ch:ru-boots-skirt-coat-shoes-socks}
\hlchaptermodality{102}

\begin{chapteropening}
\textbf{By the end of this chapter:} \emph{I can say boots, a skirt, a coat, shoes and socks.}

Complete the last lesson of chapter 102: I can say boots, a skirt, a coat, shoes and socks.
\end{chapteropening}

\section[\texorpdfstring{\ru{сапоги} (sapogí)}{сапоги (sapogí)}]{\ru{сапоги} (sapogí) — boots}

\label{lesson:RU-C102-boots}



\begin{quote}
Before the new one: say the Russian for gloves, then the Russian for a scarf.
\end{quote}

\subsection*{You'll want to know: \ru{сапоги}}
\textbf{\ru{сапоги}} (\emph{sapogí}) — \textquotedblleft{}boots\textquotedblright{}.

\begin{cousinweb}
\textbf{\ru{сапоги}} is the plural of \emph{\ru{сапог}}, \textquotedblleft{}a boot\textquotedblright{}, Old Church Slavonic \emph{sapogŭ}. Vasmer reports the usual link with \emph{\ru{сопеть}}, \textquotedblleft{}to wheeze\textquotedblright{}, through \textquotedblleft{}a pipe\textquotedblright{} and so \textquotedblleft{}a boot shaft\textquotedblright{}; Trubachev, in a note, takes it from early Bulgar or Turkic \emph{*sapag}, \textquotedblleft{}a boot with a shaft\textquotedblright{}, from a word for \textquotedblleft{}stem\textquotedblright{}.
\end{cousinweb}

\begin{grammarlens}[title={What you've built}]
One more word for this chapter.
\end{grammarlens}

\subsection*{Your turn}
\begin{itemize}
\raggedright
  \item \emph{Say it:} \emph{sapogí}
  \item \emph{Say it:} \emph{sapogí}, once more
  \item \emph{Say it:} say \emph{sapogí}
  \item \emph{Recall:} say the Russian for gloves, then the Russian for a scarf, then say \emph{sapogí} again
\end{itemize}

\begin{tcolorbox}[breakable,colback=teal!4,colframe=teal!35!black,title={Before you move on}]
What does \ru{сапоги} mean? (\textquotedblleft{}Boots\textquotedblright{}.) Say it once more.
\end{tcolorbox}
\section[\texorpdfstring{\ru{юбка} (yúbka)}{юбка (yúbka)}]{\ru{юбка} (yúbka) — a skirt}

\label{lesson:RU-C102-skirt}



\begin{quote}
Before the new one: say the Russian for a scarf, then the Russian for boots.
\end{quote}

\subsection*{You'll want to know: \ru{юбка}}
\textbf{\ru{юбка}} (\emph{yúbka}) — \textquotedblleft{}a skirt\textquotedblright{}.

\begin{cousinweb}
\textbf{\ru{юбка}} — older \emph{\ru{юпка}}, from 1701; it came through Polish \emph{jupa}, \textquotedblleft{}a jacket\textquotedblright{}, from Middle High German \emph{joppe}, which goes back through Old French \emph{jupe} to Italian \emph{giuppa} and Arabic \emph{jubba}, \textquotedblleft{}an undergarment\textquotedblright{} (Vasmer).
\end{cousinweb}

\begin{grammarlens}[title={What you've built}]
One more word for this chapter.
\end{grammarlens}

\subsection*{Your turn}
\begin{itemize}
\raggedright
  \item \emph{Say it:} \emph{yúbka}
  \item \emph{Say it:} \emph{yúbka}, once more
  \item \emph{Say it:} say \emph{yúbka}
  \item \emph{Recall:} say the Russian for a scarf, then the Russian for boots, then say \emph{yúbka} again
\end{itemize}

\begin{tcolorbox}[breakable,colback=teal!4,colframe=teal!35!black,title={Before you move on}]
What does \ru{юбка} mean? (\textquotedblleft{}A skirt\textquotedblright{}.) Say it once more.
\end{tcolorbox}
\section[\texorpdfstring{\ru{пальто} (paltó)}{пальто (paltó)}]{\ru{пальто} (paltó) — a coat}

\label{lesson:RU-C102-coat}



\begin{quote}
Before the new one: say the Russian for boots, then the Russian for a skirt.
\end{quote}

\subsection*{You'll want to know: \ru{пальто}}
\textbf{\ru{пальто}} (\emph{paltó}) — \textquotedblleft{}a coat\textquotedblright{}.

\begin{cousinweb}
\textbf{\ru{пальто}} comes from French \emph{paletot}, \textquotedblleft{}an overcoat\textquotedblright{}; Vasmer traces it further to Spanish \emph{palletoque}, \textquotedblleft{}a hooded cloak\textquotedblright{}, and Latin \emph{palla}, \textquotedblleft{}a mantle\textquotedblright{}.
\end{cousinweb}

\begin{grammarlens}[title={What you've built}]
One more word for this chapter.
\end{grammarlens}

\subsection*{Your turn}
\begin{itemize}
\raggedright
  \item \emph{Say it:} \emph{paltó}
  \item \emph{Say it:} \emph{paltó}, once more
  \item \emph{Say it:} say \emph{paltó}
  \item \emph{Recall:} say the Russian for boots, then the Russian for a skirt, then say \emph{paltó} again
\end{itemize}

\begin{tcolorbox}[breakable,colback=teal!4,colframe=teal!35!black,title={Before you move on}]
What does \ru{пальто} mean? (\textquotedblleft{}A coat\textquotedblright{}.) Say it once more.
\end{tcolorbox}
\section[\texorpdfstring{\ru{туфли} (túfli)}{туфли (túfli)}]{\ru{туфли} (túfli) — shoes}

\label{lesson:RU-C102-shoes}



\begin{quote}
Before the new one: say the Russian for a skirt, then the Russian for a coat.
\end{quote}

\subsection*{You'll want to know: \ru{туфли}}
\textbf{\ru{туфли}} (\emph{túfli}) — \textquotedblleft{}shoes\textquotedblright{}.

\begin{cousinweb}
\textbf{\ru{туфли}} — recorded from 1714; it comes from Low German \emph{tuffel}, \textquotedblleft{}a shoe\textquotedblright{}, from older \emph{pantuffel}, from Italian \emph{pantofola}, \textquotedblleft{}a slipper\textquotedblright{} (Vasmer). Its earlier source is disputed.
\end{cousinweb}

\begin{grammarlens}[title={What you've built}]
One more word for this chapter.
\end{grammarlens}

\subsection*{Your turn}
\begin{itemize}
\raggedright
  \item \emph{Say it:} \emph{túfli}
  \item \emph{Say it:} \emph{túfli}, once more
  \item \emph{Say it:} say \emph{túfli}
  \item \emph{Recall:} say the Russian for a skirt, then the Russian for a coat, then say \emph{túfli} again
\end{itemize}

\begin{tcolorbox}[breakable,colback=teal!4,colframe=teal!35!black,title={Before you move on}]
What does \ru{туфли} mean? (\textquotedblleft{}Shoes\textquotedblright{}.) Say it once more.
\end{tcolorbox}
\section[\texorpdfstring{\ru{носки} (noskí)}{носки (noskí)}]{\ru{носки} (noskí) — socks}

\label{lesson:RU-C102-socks}



\begin{quote}
Before the new one: say the Russian for a coat, then the Russian for shoes.
\end{quote}

\subsection*{You'll want to know: \ru{носки}}
\textbf{\ru{носки}} (\emph{noskí}) — \textquotedblleft{}socks\textquotedblright{}.

\begin{cousinweb}
\textbf{\ru{носки}} is the plural of \emph{\ru{носок}}, \textquotedblleft{}a sock\textquotedblright{}, built on \emph{\ru{нос}}, \textquotedblleft{}nose\textquotedblright{} (Vasmer): the tip of the foot or shoe, then the short stocking that covers it.
\end{cousinweb}

\begin{grammarlens}[title={What you've built}]
One more word for this chapter.
\end{grammarlens}

\subsection*{Your turn}
\begin{itemize}
\raggedright
  \item \emph{Say it:} \emph{noskí}
  \item \emph{Say it:} \emph{noskí}, once more
  \item \emph{Say it:} say \emph{noskí}
  \item \emph{Recall:} say the Russian for a coat, then the Russian for shoes, then say \emph{noskí} again
\end{itemize}

\begin{tcolorbox}[breakable,colback=teal!4,colframe=teal!35!black,title={Before you move on}]
What does \ru{носки} mean? (\textquotedblleft{}Socks\textquotedblright{}.) Say it once more.
\end{tcolorbox}
